% TOP_PROBLEM: {"id": 2787, "title": "Kirby Problem 2.39", "category_id": 7, "category": {"id": 7, "name": "topology", "display_name": "Topology", "description": "Properties preserved under continuous deformations.", "slug": "topology", "order_index": 7, "created_at": "2026-07-31T15:26:25.671Z"}, "status": "open"}
% FIRST_POSTED: null
% ENTRY_KIND: statement_only
% Imported from: batch11.tex; UnsolvedMath ID 2787
\documentclass[11pt]{article}
\usepackage[margin=25mm]{geometry}
\usepackage{fontspec}
\setmainfont{Latin Modern Roman}
\setsansfont{Latin Modern Sans}
\usepackage{amsmath,amssymb,mathtools}
\usepackage{unicode-math}
\setmathfont{Latin Modern Math}
\usepackage{xurl,hyperref}
\setcounter{secnumdepth}{0}
\setlength{\parindent}{0pt}
\setlength{\parskip}{5pt}
\setlength{\emergencystretch}{3em}
\newcommand{\sref}[2]{\hyperlink{source:#1:#2}{[S#2]}}
\newcommand{\cataloguescope}[1]{\paragraph{Recorded target.}#1}
\begin{document}
% UnsolvedMath ID 2787; source catalogue code KP-2.39
\setcounter{section}{2786}
\section{Kirby Problem 2.39}\label{2787}
{\small\sffamily UnsolvedMath ID 2787 · Topology}
\subsection{Definitions and mathematical statement}

\paragraph{Shared notation.}$\mathbb N=\{1,2,\ldots\}$, and $\mathbb Z,\mathbb Q,\mathbb R,\mathbb C$ denote the integers, rationals, reals and complex numbers. Any other conventions are specified locally.


An oriented surface $S$ is of infinite type if it is not a compact finite-genus surface with finitely many punctures and boundary circles. Its mapping class group $\operatorname{Mod}(S)$ consists of orientation-preserving homeomorphisms modulo isotopy, fixing boundary pointwise. An end is an element of the inverse limit of the sets of components of complements of compact subsets. An end is planar if it has a genus-zero neighborhood; otherwise it is accumulated by genus. The one-ended infinite-genus surface is the boundaryless Loch Ness monster surface. A subgroup has finite index if it has finitely many cosets. On a chosen conformal surface $\Sigma$ homeomorphic to $S$, a quasiconformal homeomorphism has globally bounded dilatation; in local coordinates its Beltrami coefficient has essential supremum strictly smaller than one. Let $\operatorname{QCMod}(\Sigma)\le\operatorname{Mod}(S)$ consist of classes with quasiconformal representatives for this conformal structure. This subgroup can depend on $\Sigma$.
\paragraph{Statement recorded in the supplied source.}
Ask the following separately:
\begin{enumerate}
\item For every infinite-genus surface with no planar ends, does its mapping class group contain an isomorphic copy of every countable group?
\item Does the mapping class group of the one-ended infinite-genus surface have a proper finite-index subgroup?
\item For an infinite-type conformal surface $\Sigma$, characterize $\operatorname{QCMod}(\Sigma)$ and determine when it is normal in $\operatorname{Mod}(S)$.
\end{enumerate} \sref{2787}{2}

\subsection{Short English statement}
Ask whether all countable groups embed in the stated mapping class groups, whether the Loch Ness monster group has a proper finite-index subgroup, and how subgroups of classes admitting quasiconformal representatives are characterized and when they are normal.
\subsection{Sources}
\begin{description}
\item[\hypertarget{source:2787:1}{S1}] ULAM.AI and the UnsolvedMath Contributors, UnsolvedMath: A Curated Collection of Open Mathematics Problems, record 2787 (KP-2.39). CC BY 4.0. \url{https://huggingface.co/datasets/ulamai/UnsolvedMath}.
\item[\hypertarget{source:2787:2}{S2}] R. İnanç Baykur, Robion C. Kirby and Daniel Ruberman, eds., \emph{K3: A New Problem List in Low-Dimensional Topology}, Mathematical Surveys and Monographs 295, American Mathematical Society (2026), Problem 2.39 and its remarks; Chapters 1 and 2 for the stated conventions. \url{https://math.berkeley.edu/sites/default/files/surv-295-ruberman-watermarked-author-pdf.pdf}.
\end{description}
\label{end:2787}
\end{document}
